\documentclass{standalone}
\usepackage{circuitikz}
\usetikzlibrary{calc}
\definecolor{circuitnotemark}{HTML}{FE00FE}
\begin{document}
\begin{circuitikz}[american, line width=0.8pt]
\ctikzset{bipoles/length=1.2cm}
\ctikzset{grounds/scale=1.36}
\coordinate (a1) at (0,0);
\coordinate (c1) at (0,-4);
\coordinate (a3) at (4,0);
\coordinate (c3) at (4,-4);
\draw (a1) to[battery1, l_=$B_{1}$, a^=$3\,\mathrm{V}$] (c1); % line 3
\draw (a1) to[R, l_=$R_{1}$, a^=$330\,\Omega$] (a3); % line 4
\draw (a3) to[leD, l_=$D_{1}$] (c3); % line 5
\node[ground] at (c1) {}; % line 6
\draw (c1) -- (c3); % line 8
\node[circ] at (c1) {};
\node[anchor=south west, inner sep=2pt, circuitnotemark, font=\LARGE\bfseries] (circuittitle) at (current bounding box.north west) {X};
\path (circuittitle.south west) rectangle ++(6.614,0.853);
\end{circuitikz}
\end{document}
